\section{Micro Controller Units}
MCUs are designed for application-specific tasks, often with real-time
computing constraints. They are typically used in embedded systems, where they
control specific functions within a larger system. MCUs are characterized by
their low power consumption, small size, and integrated peripherals, making
them ideal for applications such as home automation, wearable devices, and
industrial control systems.
\begin{tcolorbox}
    \textbf{Examples}
    \begin{itemize}
        \item \textbf{Arduino}: A popular open-source platform based on easy-to-use hardware and software. It is widely used for prototyping and educational purposes.
        \item \textbf{Raspberry Pi Pico}: A low-cost microcontroller board based on the RP2040 chip, designed for embedded applications.
        \item \textbf{ESP8266/ESP32}: Wi-Fi-enabled microcontrollers that are commonly used in IoT projects for their wireless communication capabilities.
        \item \textbf{BLE Microcontrollers}: Such as the Nordic Semiconductor nRF52 series, which are designed for low-power Bluetooth applications.
        \item \textbf{NXP LCP1102/1104}: 32-bit microcontrollers designed for low-power applications, 32KB flash memory, 8KB RAM, and various peripherals for IoT applications.
        \item \textbf{Microchip PIC18F24Q10}: A 16-bit microcontroller with integrated USB and CAN interfaces, suitable for automotive and industrial applications.
    \end{itemize}
\end{tcolorbox}
\subsection{MPU vs MCU}
MPUs (Microprocessor Units) and MCUs (Microcontroller Units) are both types of
integrated circuits used in computing, but they serve different purposes and
have distinct characteristics.
\begin{tcolorbox}
    \textbf{MPU (Microprocessor Unit)}
    \begin{itemize}
        \item Designed for general-purpose computing tasks.
        \item Typically requires external components such as memory and peripherals to
              function.
        \item Used in applications like personal computers, servers, and high-performance
              embedded systems.
        \item Examples include Intel x86 processors and ARM Cortex-A series.
    \end{itemize}
\end{tcolorbox}
\begin{tcolorbox}
    \textbf{MCU (Microcontroller Unit)}
    \begin{itemize}
        \item Designed for specific control-oriented tasks.
        \item Integrates memory, peripherals, and processing capabilities on a single chip.
        \item Used in embedded systems, IoT devices, and applications requiring low power
              consumption.
        \item Examples include Arduino, Raspberry Pi Pico, and ESP8266/ESP32.
    \end{itemize}
\end{tcolorbox}
What about ARM? ARM isn't a company that produces microcontrollers, but rather a company that designs the architecture and licenses it to other manufacturers. Many microcontrollers, including those used in IoT devices, are based on ARM architecture. For example, the ARM Cortex-M series is widely used in microcontrollers for embedded applications. So while ARM itself doesn't produce microcontrollers, its architecture is fundamental to many of the MCUs available on the market.
What is the difference between Power and Energy? Power is istantaneous, while energy is the total amount of work done over time, in fact we can calculate energy as: $E(t) = \int_{t_0}^{t} P(t) dt$.
what is a duty cycle? Duty cycle is the percentage of time a system is active compared to the total time it could be active.
